\medskip
\hypertarget{native-108-article079}{}\pdfbookmark[2]{79.}{article-079}
\begin{center}79.\end{center}
Completely determining a conic section in its plane requires \emph{three} quantities, namely the position of perihelion, the eccentricity, and the semiparameter. If these are to be derived from given quantities that depend on them, there must be enough data to form three mutually independent equations. Each radius vector given in magnitude and position supplies one equation. Determining the orbit therefore requires three radius vectors given in magnitude and position. If there are only two, either one element must itself already be given, or at least some other quantity from which a third equation can be formed. This gives rise to various problems, which we shall now consider in turn.
\par
Let $r$, $r'$ be two radius vectors that make angles $N$, $N'$ with an arbitrary straight line drawn from the Sun in the plane of the orbit, the angles being reckoned in the direction of motion. Let \SourceLeaf{109}$\Pi$ be the angle that the radius vector at perihelion makes with the same line, so that the true anomalies corresponding to the radius vectors $r$, $r'$ are $N-\Pi$, $N'-\Pi$. Finally, let $e$ be the eccentricity and $p$ the semiparameter. We then have the equations
\[
\begin{aligned}
\frac{p}{r}&=1+e\cos(N-\Pi)\\
\frac{p}{r'}&=1+e\cos(N'-\Pi),
\end{aligned}
\]
from which the other two quantities can be determined if one of $p$, $e$, $\Pi$ is given as well.
\par
Suppose first that the semiparameter $p$ is given. The quantities $e$ and $\Pi$ can clearly be determined from the equations
\[
\begin{aligned}
e\cos(N-\Pi)&=\frac{p}{r}-1\\
e\cos(N'-\Pi)&=\frac{p}{r'}-1
\end{aligned}
\]
by applying Lemma III of the preceding article. We thus have
\[
\begin{aligned}
\operatorname{tang}(N-\Pi)&=\operatorname{cotang}(N'-N)-\frac{r(p-r')}{r'(p-r)\sin(N'-N)}\\[5pt]
\operatorname{tang}(\tfrac12N+\tfrac12N'-\Pi)&=\frac{(r'-r)\operatorname{cotang}\tfrac12(N'-N)}{r'+r-\dfrac{2rr'}{p}}\mathpunct{\cdot}
\end{aligned}
\]
